\documentclass[10pt,english]{article}

\usepackage{amsmath,amssymb,amsthm}
%\usepackage{enumerate}
%\usepackage{stfloats}
\usepackage{bbm}
%\usepackage{framed}

\usepackage{tikz,pgf,pgfplots}
\usepgflibrary{shapes}
\usetikzlibrary{%
  arrows,%
  arrows.meta,
  backgrounds,
  shapes.misc,% wg. rounded rectangle
  shapes.arrows,%
  shapes,%
  calc,%
  chains,%
  matrix,%
  positioning,% wg. " of "
  scopes,%
  decorations.pathmorphing,% /pgf/decoration/random steps | erste Graphik
  shadows,%
  backgrounds,%
  fit,%
  petri,%
  quotes
}

\pgfplotsset{compat=1.12}

%%  Dimensions
%%
\topmargin             0in
\headheight            0in
\headsep               0in
\textheight          8.5in
\textwidth           6.5in
\evensidemargin        0in
\oddsidemargin         0in

\newcommand{\ul}[1]{\underline{#1}}
\renewcommand{\Pr}{\mathbb{P}}

\newcommand{\mc}[1]{\mathcal{#1}}
\newcommand{\mbb}[1]{\mathbb{#1}}
%\newcommand{\expt}{\mbb{E}}
%\newcommand{\dd}{\mathrm{d}}
\input{macros}

\begin{document}



\title{ECE 586: Vector Space Methods \\ Lecture Outlines}
\author{Henry D. Pfister \\ \univ}
\date{}
%\date{August 20th, 2020}
\maketitle

\section*{Lecture 1}

\begin{itemize}
  \item Introduce yourself
  \item Review Video-0 (intro) slides focusing on questions.  They should have watched the video.
  \item Review Video-1 slides briefly.  Focus on slides 4-6 (i.e., conditional connective and meta statement).  They should have watched the video.
  \item Have them complete in class one part of problems 2 and 3 from homework 1.
  \item Do one truth table example?
\end{itemize}

\section*{Lecture 2}

While very briefly reviewing slides,  emphasize:
\begin{itemize}
  \item Focus on slides 3-4 with multiple quantifiers.
  \item Prove ``$\sqrt{2}$ is irrational'' (see slide 6).  First, start with this weak formulation and walk through the formalization to ``If $x \in \mathbb{R}$ and $x^2=2$, then $x$ is not rational'' while defining all terms along the way.
  \begin{itemize}
    \item Suppose $x$ is a non-negative rational number and $x^2=2$ where the rational numbers are defined by $\mathbb{Q} = \{ p/q \in \mathbb{R} \,|\, p\in \mathbb{Z}, q \in \mathbb{N} \}$.
    \item Since reducing a fraction does not change its value, we can further assume that $p/q$ in reduced form where $p$ and $q$ are coprime integers (i.e., they have no common factor greater than 1).
    \item This implies that there exist integers $p>0$ and $q>0$ such that $(p/q)^2 = 2$ or $p^2 = 2 q^2$.
    \item Thus, $p^2$ is an even number (note: one should define even formally as $\{e \in \mathbb{Z} \,|\, \exists k \in \mathbb{Z}, e=2k \}$).
    \item This implies that $p$ itself is even (i.e., one can prove that '' if $k$ is not even, then $k^2$ is not even).
    \item Because $p$ is even, there exists an integer $r$ satisfying $p = 2r$.
    \item  We then obtain the equation $(2r)^2 = 2q^2$, which is equivalent to $2r^2 = q^2$ after simplification.
    \item  Because $2r^2$ is even, it follows that $q^2$ is even, which means that $q$ is also even.
    \item We conclude that $p$ and $q$ are both even. But this contradicts the fact that $p/q$ is irreducible.
    \item Hence, there is no rational $x$  such that $x^2 = 2$ and, thus  $\sqrt{2}$ is irrational.
  \end{itemize}
  \item Prove sum identity from slide 6 via induction.  Emphasize that one proves $P(n) \to P(n+1)$ by writing down $P(n+1)$ and then using $P(n)$ to simplify.
\end{itemize}

\section*{Lecture 3}

While very briefly reviewing slides,  emphasize:
\begin{itemize}
  \item Talk about countability of rationals and $\mathbb{N}^2$ and $\mathbb{Z}^2$.  For any fixed element in set, enumeration must show element in finite time.  What orderings don't work?
  \item Emphasize connection between set theory and logic
  \item Show logical De Morgan's is equivalent to set theoretical De Morgan's
  \item Emphasize definition of infinite union and intersection does not require sequences or limits
  \item Describe Russell's paradox and note problem is unrestricted comprehension
  \item In this class, we avoid issues by working with explicitly defined sets
\end{itemize}
Add homework problems


\section*{Lecture 4}

While very briefly reviewing slides,  emphasize:
\begin{itemize}
  \item Give examples of $\leq,\geq,<,>,=$ as relations
  \item Discuss birthday example along with equivalence classes
  \item State powerset example where $A = \mathcal{P}(\{1,2,3\})$ and $X,Y \in A$ are equivalent (i.e., $X\sim Y$) if $|X|=|Y|$.
  Ask students to find equivalence classes and quotient set for this case. 
  \item For functions, start with the thought question, is $x^2$ invertible? (A: depends on domain!)
  \item $f:\mathbb{R} \to \mathbb{R}$, $f(x)=x^2$ is not invertible but restriction to $\mathbb{R}_{\geq 0}$ is.
  What about $x^2 + 1$?
\end{itemize}

\section*{Lecture 5}

While very briefly reviewing slides,  emphasize:
\begin{itemize}
  \item Foundational role of real numbers with absolute distance
  \item Prove $|a+b| \leq |a|+ |b|$ using two cases: $a,b$ non-zero with same sign and its negation
  \item For space of continuous functions, emphasize above inequality and separate optimization
  \item From $x_1 = 2$, show rationals $x_{n+1} = \frac{x_n}{2} + \frac{1}{x_n}$ satisfy $x_n - \sqrt{2} \to 0$
  \item Topology Q1: ``open ball'' is open according to defn? (Write formal  logical statement, don't solve)
  \item Topology Q2: Show union of two open sets is open using defn of open
  \item Draw pictures of interior, limit point, closure, and boundary   
\end{itemize}


\section*{Lecture 6}

While very briefly reviewing slides,  emphasize:
\begin{itemize}
  \item Extended real numbers as metric space and topology of this space
  \item Definition of supremum: ``upper bound'' and ``no smaller upper bound''
  \item Prove that the definition of supremum implies a sequence achieving supremum
  \item Use this to prove a bounded non-decreasing real sequence converges to supremum
  \item Apply this to show that, if a sum is absolutely convergent, then the sum converges
  \item Draw pictures and use examples to explain different types of continuity 
\end{itemize}


\section*{Lecture 7}

While very briefly reviewing slides,  emphasize:
\begin{itemize}
  \item completeness as a desirable property that you can choose when setting up the problem
  
  \item Student question: show a closed subset of a complete space is complete

  \item L2 function distance integral?

  \item Explain contraction pictures

  \item Prove contraction mapping theorem
  
  \item Explain $f(x) = \cos(x) $ example
  
  \item Highlight useful applications
\end{itemize}  

\section*{Lecture 8}

define complete and compact for subsets

closed subset of complete space is complete

closed subset of a compact space is compact

Prove Bolzano-Weierstrauss for reals

Compactness subsequence proof

Discuss applications to probability?

continuous function on a compact set

\section*{Lecture 9}

emphasize different views of matrix product

class question: Matrix for elementary row op

discuss invertible

class question: Why is product of invertible invertible?

foundation of dimension
Class question: Why?


  
%\bibliographystyle{alpha}
%\footnotesize
%\bibliography{IEEEabrv,WCLabrv,WCLbib,WCLnewbib}
%\end{frame}


\end{document}
